% sections/discussion.tex -- the conclusion (paper_v2). Until B4 (25 November) this file held
% Sections 8-10 (Discussion, Threats to validity, Conclusion), cut in C6 from the full version's. B4
% brought in the reviewed Vietnamese version (paper_v2_vn, L10, N26 and N30): why the searched sets
% lose and why GAT leads are in Section 7 (7.4 and 7.5), and this conclusion holds the verdict that
% depends on the context, the lesson on who should judge a search, the limitations (with the two
% minor ones of the old Threats list) and future work. The other Threats items are in Section 6 or
% in the supplementary document. Claims stay within the full version's, and every number is \val{}
% from numbers.tex.
\section{Conclusion}
\label{sec:conclusion}

On a cleaned Amazon-Meta with three node types and \val{search-cands} candidate paths, automatic
meta-path search comes fairly close to manual design for HAN-based product recommendation: the sets
that a \val{search-runs-total}-run search chooses trail the manual sets by
\val{han-6-hit-relabs}--\val{han-4-hit-relabs}\% on the metric fixed in advance and tie with them on
full ranking. The loss sits mostly on the products that the searched sets' category paths cannot
reach. Which way to choose therefore depends on the context. If an expert can design the paths in
less time than the search takes (about \val{search-total-h} GPU hours here), manual design is the
better choice. Without an expert, automatic search gives sets close to manual design with no domain
knowledge. However, the random control shows that the choice of paths matters little on this small
schema: random draws from the same candidates cost nothing and do at least as well as both ways
(Section~\ref{sec:res-random}). The advantage of search (if any) more likely lies in larger schemas
that are hard both to design by hand and to sample at random, which this study does not test. A
type-blind GAT beats both meta-path models.

The main lesson is that a search for a recommender should be judged by that recommender (as RMS does
\citep{ning2022automatic}) or at least checked against random draws from its own candidates. The
sets that LMSPS's own network judged best are not HAN's best. LMSPS's gains in HAN
\citep[Table~4]{li2024longrange} were measured on node classification, where every node is encoded
from its neighbourhood. In our link prediction, most of HAN's parameters are free user vectors that
no meta-path shapes.

\paragraph{Limitations.} The study uses one dataset, one task and one small schema, with
product-side paths only. Its single random split ignores review dates, which lets a model learn from
later interactions and can change how models rank \citep{ji2023leakage,meng2020splitting}. Every
set uses one untuned configuration. With a shorter path-weight warm-up, the searched set beats the
manual one at six paths in a check with a single seed (Section~\ref{sec:res-sensitivity}). Our
search is smaller than the original paper's because it was sized for one GPU. Its runs agree beyond
chance, but a larger search might choose differently. The manual sets were designed from stated
principles rather than by a domain expert, so an expert might choose other sets. HAN's training had
not fully converged: validation BPR still improved by up to \val{han-gain-max} in the last ten
epochs, but by at most \val{searched-gain-max} for the searched sets. Longer training would
therefore more likely widen than close the gap between the two ways. Products link only to their
leaf categories (Section~\ref{sec:data}). With every ancestor linked, the candidates and the
search's choices would differ.

Future work includes (1) whether a dense category path closes the gap, (2) a search judged by its
recommender (as in RMS), (3) user-side meta-paths, (4) the top-$k$ query of \citet{zhu2020topk}
and (5) larger schemas, other datasets and splits by date, where a search may pay off.
